\section{Formalization}\label{sec:formal}

This section presents the our choice logic and choice types. Similar with separation logic, our choice logic is also based on underline resource algebra, although such an algebra is additionally equipped with the branching operators. In the rest of this section, we will first introduce the preliminary definitions and then present the choice algebra, choice logic, and our choice types.

\subsection{Preliminaries}

In functional language, the context of a program typically be treated as a store or substitution. We provides the corresponding definition and related operations as the following.

\begin{definition}[Substitution]\label{def:subst-compat}
  Substitutions are partial maps from variables to values, i.e., $\Code{Var} \rightharpoonup \Code{Value}$. We write $\Dom{\subs}$ for its domain, i.e., the set of variables where the substitution is defined.
  Write $\mapCompat{\sigma_1}{\sigma_2}$ when $\sigma_1$ and $\sigma_2$ are
  \emph{compatible}, i.e.\ they agree on the intersection of their domains.
  \[
    \mapCompat{\sigma_1}{\sigma_2}
    \doteq
    \forall x \in \Dom{\sigma_1} \cap \Dom{\sigma_2}.\; \sigma_1(x) = \sigma_2(x).
  \]
  Under $\mapCompat{\sigma_1}{\sigma_2}$, the \emph{union} of substitutions is again a substitution, written $\sigma_1 \cup \sigma_2$ as usual. Fix finite $X \subseteq_{\mathrm{fin}} \Code{Var}$ and a substitution $\sigma$, we write $\restrictTo{\sigma}{X}$ for the restriction of the $\sigma$ to $\Dom{\sigma} \cap X$, i.e.,
  \begin{align*}
    \restrictTo{\sigma}{X} \doteq \{ [x\mapsto \sigma(x)] \mid x \in \Dom{\sigma} \cap X \}.
  \end{align*}
 \end{definition}

 As mentioned in \autoref{sec:overview}, the choice resource is an capablity to freely decide assignments of variables. Therefore, we define the choice resource the following.

\begin{definition}[Resource]\label{def:resource}
A resource is a \emph{non-empty} set of substitutions that \emph{share the same domain}. The resource space is defined as the following.
\begin{align*}
  \Res \doteq \{ m ~|~ m \neq \emptyset \text{ and } m \in \mathcal{P}(\Code{Var} \rightharpoonup \Code{Value}) \text{ and } \forall \sigma\;\sigma' \in m. \Dom{\sigma} = \Dom{\sigma'} \}.
\end{align*}
We also extends the definition of domain, compatible, and restriction from substitutions to resources. That is,
\begin{align*}
  &\Dom{m} \doteq  \Dom{\subs} \quad \text{for arbitrary } \subs \in m 
  \\&\mapCompat{m_1}{m_2}
  \doteq
  \forall \subs_1 \in m_1.\;\forall \subs_2 \in m_2.\; \mapCompat{\subs_1}{\subs_2}
  \\&\restrictTo{m}{X}
  \doteq
  \{\, \restrictTo{\subs}{X} \mid \subs \in m \,\}
\end{align*}
\end{definition}

Note that for any resource $m$, $\mapCompat{m}{m}$ always holds since they always agree on their domain. Similarly, singleton set $\{m\}$ is compatible with itself.

\begin{definition}[Resource Operation]\label{def:resource-times} We define two operations ($\Res \times \Res \rightharpoonup \Res$) for resources. The \emph{resource product} $m_1 \times m_2$ is the partial binary operation which produce peicewise union of two resources, i.e.,
  \[
    m_1 \times m_2 \;\doteq\;
    \begin{cases}
      \{\, \subs_1 \cup \subs_2 \mid \subs_1 \in m_1,\; \subs_2 \in m_2 \,\}
        & \text{if } \mapCompat{m_1}{m_2}, \\
      \text{undefined}
        & \text{otherwise.}
    \end{cases}
  \]
The \emph{resource sum} $m_1 + m_2$ is the union of two resources, i.e.,
\[
  m_1 + m_2 \;\doteq\;
  \begin{cases}
    m_1 \cup m_2
      & \text{if } \forall m_1 \in m_1.\; \forall m_2 \in m_2.\; \Dom{m_1} = \Dom{m_2}, \\
    \text{undefined}
      & \text{otherwise.}
  \end{cases}
\]
  \end{definition}

The partial order $\leq$ over resources is the analogue of ``more information by releasing the variables set constraint''.

\begin{definition}[Partial order]\label{def:partial-order} The partial order $\leq$ is defined as 
    \begin{align*}
      m_1 \leq m_2 \doteq \restrictTo{m_2}{\Dom{m_1}} = m_1.
    \end{align*}
Easily to know, $\leq$ is reflexive, antisymmetric, and transitive.
\end{definition}

\begin{definition}[Restriction-agreeing substitutions]\label{def:res-agree-in-resource}
  Given a substitution $\sigma$ and a resource $m$, we write $\resAgree{m}{\sigma}$ (suggestively, $m$ semijoin $\sigma$) for the set of substitutions in $m$ whose restriction to $\Dom{\sigma}$ equals~$\sigma$:
  \begin{align*}
    \resAgree{m}{\sigma} \doteq \{\, \sigma' \mid \sigma' \in m \text{ and } \restrictTo{\sigma'}{\Dom{\sigma}} = \sigma \,\}.
  \end{align*}
\end{definition}

\subsection{Choice algebra}

The choice resource forms an algebra, which is not only a partial commutative monoid as in separation logic, but also a partial commutative semiring (PCS).

\begin{definition}[Choice algebra]\label{def:choice-semiring}
  A \emph{partial commutative semiring without zero}
  $(\mathcal{M},{+},{\times},\mathbf{1})$ consists of a carrier set $\mathcal{M}$ and partial binary operations ${+}$ and ${\times}$ and unit elements $\mathbf{1}$ that satisfies the following properties:
  \begin{itemize}
  \item Indentity: $\forall m \in \mathcal{M}. m \times \mathbf{1} = m$.
  \item Commutative (when operations are defined): $\forall m_1\;m_2 \in \mathcal{M}. m_1 + m_2 = m_2 + m_1$ and $m_1 \times m_2 = m_2 \times m_1$.
  \item Associative (when operations are defined): $\forall m_1\;m_2\;m_3 \in \mathcal{M}. (m_1 + m_2) + m_3 = m_1 + (m_2 + m_3)$ and $(m_1 \times m_2) \times m_3 = m_1 \times (m_2 \times m_3)$.
  \end{itemize}
The choice algebra is a partial commutative semiring equipped with a partial order $\leq$ that satisfies the bifunctoriality property:
  \begin{align*}
    \forall m_1\;m_2\;m_1'\;m_2' \in \mathcal{M}. m_1 \leq m_1' \land m_2 \leq m_2' &\implies m_1 \times m_2 \leq m_1' \times m_2' 
    \\& (\text{when } m_1 \times m_2 \text{ and } m_1' \times m_2' \text{ are defined.})
    \\\forall m_1\;m_2\;m_1'\;m_2' \in \mathcal{M}. m_1 \leq m_1' \land m_2 \leq m_2' &\implies m_1 + m_2 \leq m_1' + m_2' 
    \\& (\text{when } m_1 + m_2 \text{ and } m_1' + m_2' \text{ are defined.})
  \end{align*}
  \end{definition}

Easily to know, $(\Res,{+},{\times},\{ \emptyset \})$ and $\leq$ is a choice algebra.

\begin{example}[Choice algebra]
Let $m_x \doteq \{\,[x \mapsto 1],\, [x \mapsto 2]\,\}$, so
$\Dom{m_x} = \{x\}$.
Let $m_y \doteq \{\,[y \mapsto 3],\, [y \mapsto 4]\,\}$, with $\Dom{m_y} = \{y\}$.
Then
\[
  m_x \times m_y
  = \left\{
    \begin{array}{@{}l@{}}
      [x \mapsto 1,\, y \mapsto 3],\,
      [x \mapsto 1,\, y \mapsto 4],\\[\jot]
      [x \mapsto 2,\, y \mapsto 3],\,
      [x \mapsto 2,\, y \mapsto 4]
    \end{array}
  \right\}.
\]
Now take
$m_{xy} \doteq \{\, [x \mapsto 1,\, y \mapsto 3],\, [x \mapsto 2,\, y \mapsto 4]\,\}$.
We have
$\restrictTo{m_{xy}}{\{x\}} = m_x$ and $\restrictTo{m_{xy}}{\{y\}} = m_y$,
so by the definition of~$\leq$,
$m_x \leq m_{xy}$ and $m_y \leq m_{xy}$.
The product satisfies the same restrictions:
$\restrictTo{(m_x \times m_y)}{\{x\}} = m_x$ and
$\restrictTo{(m_x \times m_y)}{\{y\}} = m_y$.
Nevertheless $|m_x \times m_y| = 4 > 2 = |m_{xy}|$; equivalently,
\[
  m_x \times m_y
  = m_{xy} + \{\, [x \mapsto 1,\, y \mapsto 4],\, [x \mapsto 2,\, y \mapsto 3] \,\},
\]
\end{example}

\paragraph{Partial order, addition, and multiplication} We always have
\begin{align*}
  \forall m_1\;m_2 \in \Res. m_1 \leq m_1 \times m_2 \quad (\text{when } m_1 \times m_2 \text{ is defined.})
\end{align*}
that is because we always have $\restrictTo{m_1 \times m_2}{\Dom{m_1}} = m_1$. However, the similar property does not hold for addition: when $m_1 + m_2$ is defined we do \emph{not}
generally expect the usual monotonicity inequality
$m_1 \leq m_1 + m_2$.
This is unlike incorrectness logic~\cite{OHP19} and outcome
logic~\cite{ZDS23,Zil25TOPLAS}, where branching is tightly coupled to
the consequence relation.
Here, $+$ simply adjoins more substitutions---i.e.\ it records strictly
more alternatives---and whether that strengthens or weakens a
judgment depends on the approximation mode.
Under an overapproximate interpretation, extra alternatives relax universal
guarantees (more ways a safety property can fail); under an
underapproximate interpretation, they add witnesses for existential or
reachability claims.
Because the two modes are dual and our modalities can mark formulas
with either interpretation, a logic that uniformly enforces Kripke monotonicity
should not identify $\leq$ with enrichment by~$+$.
By contrast, expanding or shrinking the set of tracked variables is
orthogonal to over- versus under-approximation: for example, the same
resource
$\{\,[x \mapsto 0],\, [x \mapsto 1],\, [x \mapsto 2]\,\}$
is compatible with an overapproximate interpretation that $x$ may range
over~$[0,5]$ and, separately, with an underapproximate interpretation that
witnesses for~$x$ lie in~$[1,2]$.
Extending every substitution with arbitrary bindings for a fresh
variable~$y$ disturbs neither kind of judgment.

\subsection{Choice logic}

\emph{Basic} choice logic is parameterized by set of atomic propositions $\mathcal{A}$ and has no approximation or persistency modality. Then, every judgment is \emph{precisely} describes a resource.
We write $m \models p$ for satisfaction of a resource proposition $p$
at world $m \in \Res$.

\begin{definition}[Basic choice logic] Let $(\mathcal{M},{+},{\times},\mathbf{1}, \leq)$ to be a choice algebra. Let $\mathcal{A}$ to be atomic propositions. Let $\denot{-} : \mathcal{A} \sarr \mathcal{M}$ be an interpretation function of pure formulas $\phi$. Then we have:
\begin{itemize}
\item $m \models \chtrue$ always;
\item $m \models \chfalse$ never;
\item $m \models a$ iff $a \in \mathcal{A}$ and $\denot{a} \leq m$;
\item $m \models p \chland q$ iff $m \models p$ and $m \models q$;
\item $m \models p \chlor q$ iff $m \models p$ or $m \models q$;
\item $m \models p \chimpl q$ iff for all $m \leq m'. m' \models p$ implies $m' \models q$;
% \item $m \models \forall x.\,p$ iff for all values $v$, $m \models p[x\mapsto v]$;
% \item $m \models \exists x.\,p$ iff for some $v$, $m \models p[x\mapsto v]$.
\item $m \models p \chstar q$ iff $\exists m_1\;m_2.\; m_1 \times m_2 \leq m$ and
  $m_1 \models p$ and $m_2 \models q$.
\item $m \models p \chwand q$ iff $\forall m'.\; \mapCompat{m'}{m}$ and $m' \models p$ implies
  $m' \times m \models q$.
\item $m \models p \choplus q$ iff $\exists m_1\;m_2.\; m_1 + m_2 \leq m$ and
  $m_1 \models p$ and $m_2 \models q$.
\end{itemize}
We still write $p \models q$ iff for all $m \in \Res$, $m \models p$ implies $m \models q$.
\end{definition}

Intuitively, $m$ models $\subExt p$ (or $\supExt p$) iff there exists a supperset (or subset) $m'$ of $m$ that models $p$. On the other hand, $m$ models $\chpers p$ iff $m$ is a empty resource or a singleton resource that models $p$.


\begin{definition}[Choice logic]\label{def:choice-logic-modalities} Choice logic additional equiped with three modalities than basic choice logic, i.e., the overapproximation modality $\subExt$, the underapproximation modality $\supExt$, and the persistency modality $\chpers$.
  \begin{itemize}
    \item $m \models \subExt p$ iff $\exists m'.\; m \subseteq m'$ and
      $m' \models p$;
    \item $m \models \supExt p$ iff $\exists m'.\; m' \subseteq m$ and
      $m' \models p$,
    \item $m \models \readM{X}{p_1} p_2$ iff $\forall \sigma. \{\sigma\} \in \denot{\subExt p_1} \land \Dom{\sigma} = X \impl \resAgree{m}{\sigma} \models \sigma(p_2) $.
    \item $m \models \chpers p$ iff $|m| \leq 1$ and
      $m \models p$,
    \end{itemize}
\end{definition}


\begin{example}[Dependent Underapproximation] There are three different underapproximations:
  \begin{enumerate}
    \item (dependent) $\supExt (0 < x) \chland \supExt (y < 3)$
    \item (independent, square) $\supExt (0 < x) \chstar \supExt (y < 3)$
    \item (dependent, triangle) $\supExt (0 < x) \chland \supExt (x \leq y < 3)\judgfail$. The actually preresentation is:
    \begin{align*}
      \supExt (0 < x) \chland \readM{\{x\}}{(0 < x)} \supExt (x \leq y < 3)
    \end{align*}
  \end{enumerate}
We have:
\begin{align*}
  \{[x=1; y=1], [x=1; y=2]\} &\models \readM{\{x\}}{(0 < x)} \supExt (x \leq y < 3) 
  \\\{[x=2; y=2]\} &\models \readM{\{x\}}{(0 < x)} \supExt (x \leq y < 3)
  \\\{[x=1; y=1], [x=2; y=2]\} &\models \readM{\{x\}}{(0 < x)} \supExt (x \leq y < 3) \judgfail 
  \\&\text{If contains $[x=1]$, must includes all possible $y$.}
  \\\{[x=1; y=1], [x=1; y=2], [x=2; y=2]\} &\models \readM{\{x\}}{(0 < x)} \supExt (x \leq y < 3) \judgok 
  \\&\text{Allow to contains multiple $x$.}
\end{align*}
The $\readM{\{x\}}{(0 < x)}$ works as a ``read'' permission for $x$, which reads some assignment of $x$ that satisfies $0 < x$. Note that it is different from
\begin{align*}
  &m \models \subExt (0 < x) \chland \supExt (x \leq y < 3) 
  \\\iff& m \models \subExt (0 < x) \text{ and } m \models \supExt (x \leq y < 3)
  \\\implies& \forall \sigma \in m. 0 < \sigma(x) \text{ and } m \models \supExt (x \leq y < 3)
  \\\implies& \forall \sigma \in m. 0 < \sigma(x) \text{ and } [x=-1;y=2] \in m \quad \judgfail
\end{align*}\noindent
The problem is that the $\chland$ cannot make the constraint $0 < x$ be added into the pure proposition $x \leq y < 3$ (i.e., to be $\supExt (0 < x \land x \leq y < 3)$).
\end{example}

When the expressiveness of atomic propositions is rich enough, the nested uses of $\subExt$ and $\supExt$ over logical connectives collapse to $\subExt$ and $\supExt$ applied to the underlying pure propositions directly. For example, $\subExt (\subExt \phi_1 \chland \subExt \phi_2)$ is equivalent to $\subExt \phi_1 \chland \subExt \phi_2$, and $\supExt (\subExt \phi_1 \chland \subExt \phi_2)$ is equivalent to $\supExt (\phi_1 \cap \phi_2)$.

\begin{theorem}[Collapse of nested modalities] Assume $\mathcal{A}$ is rich enough, i.e.,
\begin{align*}
  % &\forall X. \exists \phi \in \mathcal{A}. \denot{\phi} = \{ \subs \mid \Dom{\subs} = X \}
  % \\
  &\forall \phi_1\;\phi_2 \in \mathcal{A}. \exists \phi_3 \in \mathcal{A}. \denot{\phi_3} = \denot{\phi_1} \cap \denot{\phi_2}
  \\& \forall \phi_1\;\phi_2 \in \mathcal{A}. \exists \phi_3 \in \mathcal{A}. \denot{\phi_3} = \denot{\phi_1} \cup \denot{\phi_2}
  \\& \forall \phi_1\;\phi_2 \in \mathcal{A}. \exists \phi_3 \in \mathcal{A}. \denot{\phi_3} = \denot{\phi_1} \times \denot{\phi_2}
\end{align*}
Then, all formula in choice logic can be expressed in approximation normal form, i.e., $\subExt$ and $\supExt$ are only applied to pure propositions.
\end{theorem}

\subsection{Core Language and Choice Types}

\begin{figure}[h]
{\small
  \begin{alignat*}{2}
    \text{\textbf{Variables }}& \quad &\quad& x, y, z, \vnu, ...
    \\\text{\textbf{Data constructors }}& \quad &d ::= \quad & () ~|~ \Code{true} ~|~ \Code{false} ~|~ \Code{O} ~|~ \Code{S} ~|~ \Code{Cons} ~|~ \Code{Nil} ~|~ \Code{Leaf} ~|~ \Code{Node}
    \\\text{\textbf{Base Types}}& \quad & b  ::= \quad &  \Code{unit} ~|~ \Code{bool} ~|~ \Code{nat} ~|~ \Code{int} ~|~ ...
    \\\text{\textbf{Basic Types}}& \quad & s  ::= \quad &  b ~|~ s\sarr s
    \\\text{\textbf{Operations}}& \quad & \primop ::= \quad & {+} ~|~ {-} ~|~ {==} ~|~ {<} ~|~ {\leq} ~|~ ...
    \\ \text{\textbf{Constants }}& \quad &c ::= \quad & () ~|~ \mathbb{B} ~|~ \mathbb{Z} ~|~ d\ \overline{c}\ ...
    \\ \text{\textbf{Values }}& \quad  & v ::= \quad & c ~|~ x ~|~ \zlam{x}{s}{e} ~|~ \zfix{f}{s}{x}{s}{e}
    \\ \text{\textbf{Expressions}}& \quad & e ::=\quad & v ~|~ \err ~|~ \zlet{x}{e}{e} ~|~ \zlet{x{:}s}{e}{e} ~|~ \zlet{x{:}s}{\primop\ \overline{v}}{e} ~|~ 
    \\ & & \quad & |~  \zlet{x{:}s}{v\;v}{e} ~|~ \match{v} \overline{d\ \overline{y} \to e} 
  \end{alignat*}}
\caption{Core Language Syntax. }
\label{fig:core-lang-syntax}
\end{figure}

\begin{figure}[h]
  {\small
    \begin{alignat*}{2}
      \text{\textbf{Qualifiers}}& \quad & \phi ::= \quad & c ~|~ x ~|~ \primop\;\overline{v} ~|~ \bot ~|~ \top  ~|~ \neg \phi ~|~ \phi \land \phi ~|~ \phi \lor \phi ~|~ \phi \times \phi ~|~ \phi \impl \phi ~|~ \forall x{:}b.\, \phi ~|~ \exists x{:}b.\, \phi
      \\ \text{\textbf{Choice Types}}& \quad & \tau ::= \quad & \ort{b}{\phi} ~|~ \urt{b}{\phi} ~|~ \tau \interty \tau ~|~ \tau \unionty \tau ~|~ \tau \choplus \tau ~|~ x{:}\tau \sarr \tau ~|~ x{:}\tau \chwand \tau ~|~ x{:}b \garr \tau
      \\ \text{\textbf{Type Contexts}}& \quad & \Gamma ::= \quad & \emptyset ~|~ x{:}\tau ~|~ \Gamma \chstar \Gamma ~|~ \Gamma, \Gamma
    \end{alignat*}}
  \caption{Choice Types Syntax. }
  \label{fig:choice-types-syntax}
  \end{figure}

Here $\ort{b}{\phi}$ is an overapproximate refinement and
$\urt{b}{\phi}$ underapproximate; $\oplus$ on types mirrors resource
sum at the kind level; $\sarr$ and $\wand$ distinguish
multiplicative versus ``frame-preserving'' dependency in function
abstraction (as in the semantic clauses below); and $x{:}b \garr \tau$
is a dependent telescope familiar from refinement-type presentations.
The context former $\Gamma_1 \sep \Gamma_2$ uses the separating
combinator (syntax for $*$ at the context level), while ``,'' is the
ordinary additive conjunction of hypotheses.

\newpage
\subsection{Denotational semantics}

\paragraph{Untyped semantics} A term $e$ can be reduced into multiple values under a substitution.
\begin{align*}
  \denot{e} &: \Res \sarr \mathcal{P}(\Code{Value})
  \\\denot{e}(r) &\doteq \{ v ~|~ \sigma \in r \land \mstep{\sigma(e)}{v} \}
  \\ [x \mapsto V] &\doteq \{ [x\mapsto v] ~|~ v \in V \}
  \\ (application) &f_1\;f_2 \doteq \lambda r. \{ v ~|~ v_1 \in (f_1\;r) \text{ and } v_2 \in (f_2\;r) \text{ and } \mstep{v_1\;v_2}{v}  \}
\end{align*}

\paragraph{Verbose semantics} We add a read modality for all possible types:
\begin{align*}
  \Gamma \readT \tau &\doteq \text{for all substitution $\sigma$ in over-approximated context $\Gamma$, $\sigma(\tau)$ }
\end{align*}

\begin{align*}
  \\\denot{\tau} &: \Res \sarr (\Res \sarr \mathcal{P}(\Code{Value})) \sarr \Code{Prop_{Choice}}
  \\\denot{\Gamma \readT \tau} &\doteq \lambda r. \lambda f. \readM{\Dom{\Gamma}}{\denot{\Gamma}} (\denot{\sigma(\tau)}\;r\;f)
  \\\denot{\Gamma \readT \ort{b}{\phi}} &\doteq \lambda r.\lambda f. \forall \vnu. [\vnu \mapsto f\;r] \chimpl \readM{\Dom{\Gamma}}{\denot{\Gamma}} \subExt \phi
  \\\denot{\Gamma \readT \urt{b}{\phi}} &\doteq \lambda r.\lambda f. \forall \vnu. [\vnu \mapsto f\;r] \chimpl \readM{\Dom{\Gamma}}{\denot{\Gamma}} \supExt \phi
  \\\denot{x{:}\tau_x \sarr \tau} &\doteq \lambda r.\lambda f.\forall f_x. (\denot{\tau_x}\;r\;f_x) \chimpl \denot{\tau}\;(r \chland [\vnu \mapsto f_x\;r])\;(f\;\lambda r.\{\sigma(x) ~|~ \sigma \in r \})
  \\\denot{x{:}\tau_x \chwand \tau} &\doteq \lambda r.\lambda f.\forall f_x. (\denot{\tau_x}\;r\;f_x) \chimpl \denot{\tau}\;(r \chstar [\vnu \mapsto f_x\;r])\;(f\;\lambda r.\{\sigma(x) ~|~ \sigma \in r \})
  \\\denot{\chpers \tau} &\doteq \lambda r. \lambda f. \chpers (\denot{\tau}\;r\;f)
  \\\denot{\emptyset} &\doteq \Code{Atom}(\textbf{1})
  \\\denot{x{:}\tau} &\doteq \{ r ~|~ \denot{\tau}\;\textbf{1}\;(\lambda\_.\{ \sigma(x) ~|~ \sigma \in r \} ) \}
  \\\denot{\Gamma_1, \Gamma_2} &\doteq \denot{\Gamma_1} \chland \denot{\Gamma_2}
  \\\denot{\Gamma_1 \chstar \Gamma_2} &\doteq \denot{\Gamma_1} \chstar \denot{\Gamma_2}
\end{align*}\noindent

\begin{theorem}[Fundamental theorem]
  \begin{align*}
    \Gamma \vdash e : \tau \text{ when } \forall r \in \denot{\Gamma}. \denot{\tau}\;r\;\denot{e}
  \end{align*}\noindent
\end{theorem}

\begin{example}[Base case] Consider
\begin{align*}
  &\vdash e : \urt{\Int}{0 \leq \vnu \leq 3}
  \\\iff& \forall r \in \{\textbf{1}\}. \denot{\urt{\Int}{0 \leq \vnu \leq 3}}\; r\;\denot{e}
  \\\iff& \denot{\urt{\Int}{0 \leq \vnu \leq 3}}\; \textbf{1}\;\denot{e}
  \\\iff& \forall \vnu. [\vnu \mapsto \denot{e}(\textbf{1})] \chimpl \supExt (0 \leq \vnu \leq 3)
  \\\iff& \forall \vnu. \Code{Atom}(\{\vnu \mapsto v \mid \mstep{e}{v} \}) \chimpl \supExt (0 \leq \vnu \leq 3)
\end{align*}
The overapproximation is similar.
\end{example}

\begin{example}[Singleton Base Context] Consider
  \begin{align*}
    &r \in \denot{x{:}\urt{\Int}{0 \leq \vnu \leq 3}}
    \\\iff& \denot{\tau}\;\textbf{1}\;(\lambda\_.)
    \\\iff& \forall \vnu. [\vnu \mapsto \{ \sigma(x) ~|~ \sigma \in r \}] \chimpl \supExt (0 \leq \vnu \leq 3)
    \\\iff& [x \mapsto \{ \sigma(x) ~|~ \sigma \in r \}] \chimpl \supExt (0 \leq x \leq 3)
    \\\iff& r \chimpl \supExt (0 \leq x \leq 3)
  \end{align*}
  The overapproximation is similar.
\end{example}

\begin{example}[Nonempty context] Consider
  \begin{align*}
    &x{:}\urt{\Int}{0 \leq \vnu \leq 3}\vdash (x+1) : \urt{\Int}{1 \leq \vnu \leq 2}
    \\\iff& \forall r \in \{\denot{x{:}\urt{\Int}{0 \leq \vnu \leq 3}}\}. \denot{\urt{\Int}{1 \leq \vnu \leq 2}}\; r\;\denot{x+1}
    \\\iff& \forall r. (r \chimpl \supExt (0 \leq x \leq 3)) \chimpl \denot{\urt{\Int}{1 \leq \vnu \leq 2}}\; r\;\denot{x+1}
    \\\iff& \forall r. (r \chimpl \supExt (1 \leq x \leq 2)) \chimpl \forall \vnu. [\vnu \mapsto \denot{x+1}(r)] \chimpl \supExt (1 \leq \vnu \leq 2)
  \end{align*}
\end{example}

\begin{example}[Variable dependency] Consider
  \begin{align*}
    &x{:}\urt{\Int}{0 \leq \vnu \leq 3}\vdash \Code{range}(x, 3) : x{:}\ort{\Int}{0 \leq x \leq 3} \readT \urt{\Int}{ x \leq \vnu \leq 3}
    \\\iff& \forall r. (r \chimpl \supExt (0 \leq x \leq 3)) \chimpl \denot{x{:}\ort{\Int}{0 \leq x \leq 3} \readT \urt{\Int}{ x \leq \vnu \leq 3}}\; r\;\denot{\Code{range}(x, 3)}
    \\\iff& \forall r. (r \chimpl \supExt (0 \leq x \leq 3)) \chimpl \{x\}:\subExt (0 \leq x \leq 3) \readT \denot{\urt{\Int}{ x \leq \vnu \leq 3}}\; r\;\denot{\Code{range}(x, 3)}
    \\\iff& \forall r. (r \chimpl \supExt (0 \leq x \leq 3)) \chimpl
    \\& \{x\}:\subExt (0 \leq x \leq 3) \readT \forall \vnu. [\vnu \mapsto \denot{\Code{range}(x, 3)}(r)] \chimpl \supExt (x \leq \vnu \leq 3)
    \\\iff& \forall r. (r \chimpl \supExt (0 \leq x \leq 3)) \chimpl
    \\& \forall x. 0 \leq v_x \leq 3 \impl \forall \vnu. [\vnu \mapsto \denot{\Code{range}(x, 3)}(\resAgree{[x\mapsto v_x]}{r} )] \chimpl \supExt (v_x \leq \vnu \leq 3)
    \\& \forall x. 0 \leq v_x \leq 3 \impl \forall \vnu. [\vnu \mapsto \denot{\Code{range}(v_x, 3)}] \chimpl \supExt (v_x \leq \vnu \leq 3)
  \end{align*}
What happens if we read a free variable?
\begin{align*}
  &x{:}\urt{\Int}{0 \leq \vnu \leq 3}\vdash \Code{range}(x, 3) : x{:}\ort{\Int}{0 \leq x \leq 3} \readT \urt{\Int}{ 0 \leq \vnu \leq 3}
  \\\iff& \forall r. (r \chimpl \supExt (0 \leq x \leq 3)) \chimpl
  \\& \forall x. 0 \leq v_x \leq 3 \impl \forall \vnu. [\vnu \mapsto \denot{\Code{range}(v_x, 3)}] \chimpl \supExt (0 \leq \vnu \leq 3)
\end{align*}
whihc means that the \Code{range} function needs always cover the range $[0,3]$ when $x$ is in the range $[0,3]$. It indeed requires the \Code{range} function to be independent with $x$ (when $x$ is in $[0,3]$).
\end{example}

\paragraph{Read modality} An type context must be
\begin{align*}
  x{:}\urt{\Int}{\top}, y{:}\readM{\{x\}}{?} \urt{\Int}{0 \leq \vnu \leq x} \vdash ...
\end{align*}
A type judgement must be
\begin{align*}
  ... \vdash e : \readM{\{x\}}{?} \urt{\Int}{0 \leq \vnu \leq x}
\end{align*}
\textbf{Question:} How can I smartly know the read qualifier?
how can we smartly know the read qualifier?

\textbf{Intuition:} Keep the left-hand side resource to be minimal requirement, and use it as the hidden read qualifier.

\paragraph{Erase modality}
Write $\denot{\phi} \subseteq \Res$ for the set of resources that satisfy~$\phi$.
Then the $\subseteq$-closures above coincide with the modalities from \autoref{def:choice-logic-modalities}:
\begin{align*}
  \subExt(R) &\doteq \{\, r' \mid \exists r \in R.\; r' \subseteq r \,\},
  \\
  \supExt(R) &\doteq \{\, r' \mid \exists r \in R.\; r \subseteq r' \,\},
  \\
  \denot{\subExt p} &\doteq \subExt(\denot{p}) \\
  \denot{\supExt p} &\doteq \supExt(\denot{p}).
\end{align*}
We would like to erase the redundant proof obligations when we want to prove a underapproximation task.
Consider the following preorder on resources:
\begin{align*}
  % r \subExteq r' &\doteq \forall R. r \in \subExt(R) \iff r' \in \subExt(R)\\
  r \supExtleq r' &\doteq \forall R.\; r \in \supExt(R) \impl r' \in \supExt(R).
\end{align*}
It means that for $r \supExtleq r'$, if $r \models p$ holds, then $r' \models p$ also holds, these $r'$ are the redundant proof obligations.
Write $\supExtQ(S)$ for the set of \emph{$\supExtleq$-minimal} elements of $S \subseteq \Res$, i.e.\ those $r \in S$ with no strictly $\supExtleq$-smaller witness in~$S$:
\begin{align*}
  \supExtQ(S) &\doteq
  \{\, r \in S \mid \forall r' \in S.\; r' \supExtleq r \impl r = r' \,\}.
\end{align*}
Intuitively, $\supExt$ \emph{fattens} a set of resources (upward closure), whereas $\supExtQ$ keeps only $\supExtleq$-minimal witnesses (erase redundancy); heuristically the two sit at opposite ends of the same Galois-style pairing familiar from posets and Alexandrov spaces. For example, we have:
\begin{align*}
  &\supExtQ(R_1) \subseteq R_2 \iff R_1 \subseteq \supExt(R_2)
\end{align*}\noindent
We define the erase modality by collecting these minimal witnesses in the denotation of~$p$:
\begin{align*}
  % \denot{\subExtQ \phi} &= \supExtQ(\denot{\phi}) \\
  \denot{\supExtQ p} &\doteq \supExtQ(\denot{p}).
\end{align*}
When a property (not need to be underapproximated) appears at negative position, the erase modality can be used to erase the redundant proof obligations.
\begin{align*}
  &\forall p_1\;p_2.  \supExtQ p_1 \impl p_2 \implies p_1 \impl p_2
\end{align*}\noindent
Note that the following fact holds, where $\subExt \supExtQ p $ means ``flip'' the property $p$ (at negative position) from under to over:
\begin{align*}
  &\forall \phi\;p.  \subExt \supExtQ \supExt \phi \impl p \implies \subExt \phi \impl p
  \\&\forall \phi\;p. \subExt \supExtQ \subExt \phi \impl p \implies \subExt \phi \impl p
\end{align*}


\begin{align*}
  \\\denot{\tau} &: \Res \sarr (\Res \sarr \mathcal{P}(\Code{Value})) \sarr \Code{Prop_{Choice}}
  % \\\denot{\Gamma \readT \tau} &\doteq \lambda r. \lambda f. \readM{\Dom{\Gamma}}{\denot{\Gamma}} \denot{\sigma(\tau)}\;r\;f
  \\\denot{\ort{b}{\phi}} &\doteq \lambda r.\lambda f. \forall \vnu.  [\vnu \mapsto f\;r] \chimpl \Code{FV}(\phi) : \subExt r \readT \subExt \phi
  \\\denot{\urt{b}{\phi}} &\doteq \lambda r.\lambda f. \forall \vnu.  [\vnu \mapsto f\;r] \chimpl \Code{FV}(\phi) : \subExt r \readT \supExt \phi
  \\\denot{x{:}\tau_x \sarr \tau} &\doteq \lambda r.\lambda f.\forall f_x. (\denot{\tau_x}\;r\;f_x) \chimpl \denot{\tau}\;(r \chland [\vnu \mapsto f_x\;r])\;(f\;x)
  \\\denot{x{:}\tau_x \chwand \tau} &\doteq \lambda r.\lambda f.\forall f_x. (\denot{\tau_x}\;r\;f_x) \chimpl \denot{\tau}\;(r \chstar [\vnu \mapsto f_x\;r])\;(f\;x)
  \\\denot{\chpers \tau} &\doteq \lambda r. \lambda f. \chpers (\denot{\tau}\;r\;f)
  \\\denot{\emptyset} &\doteq \Code{Atom}(\textbf{1})
  \\\denot{x{:}\tau} &\doteq \supExtQ \{ r ~|~ \denot{\tau}\;\textbf{1}\;(\lambda\_.\{ \sigma(x) ~|~ \sigma \in r \} ) \}
  \\\denot{\Gamma_1, \Gamma_2} &\doteq \denot{\Gamma_1} \chland \denot{\Gamma_2}
  \\\denot{\Gamma_1 \chstar \Gamma_2} &\doteq \denot{\Gamma_1} \chstar \denot{\Gamma_2}
\end{align*}\noindent
